% ============================================================
% Relatório: análise e redução do volume de tráfego do FACEIT (CS2)
% GRUPO 6 - estilo do relatório CAIDA Ark (article 11pt, A4, margens 2,5 cm)
% Compilar com: pdflatex relatorio_faceit.tex  (duas vezes)
% As figuras fig_udp.png, fig_sni.png e fig_dns.png devem estar na mesma pasta.
% ============================================================
\documentclass[11pt,a4paper]{article}

\usepackage[utf8]{inputenc}
\usepackage[T1]{fontenc}
\usepackage{lmodern}
\usepackage[portuguese]{babel}
\usepackage[a4paper,margin=2.5cm]{geometry}
\usepackage{amsmath}
\usepackage{booktabs,array,tabularx,graphicx}
\usepackage{listings}
\usepackage{pgfplots}
\pgfplotsset{compat=1.16}
\usepackage{hyperref}

% Mesmos rótulos do relatório anterior (Abstract, Table, Figure)
\addto\captionsportuguese{%
  \renewcommand{\abstractname}{Abstract}%
  \renewcommand{\tablename}{Table}%
  \renewcommand{\figurename}{Figure}}

\renewcommand{\labelitemi}{-}
\emergencystretch=3em
\urlstyle{tt}
\renewcommand{\arraystretch}{1.15}
\newcolumntype{Y}{>{\raggedright\arraybackslash}X}

\lstset{
  basicstyle=\ttfamily\footnotesize,
  breaklines=true,
  frame=single,
  columns=fullflexible,
  keepspaces=true,
  showstringspaces=false
}

% Captura de ecrã: usa o ficheiro se existir, senão mostra uma caixa
\newcommand{\figecra}[1]{%
  \IfFileExists{#1}%
    {\includegraphics[width=\linewidth]{#1}}%
    {\fbox{\parbox{0.92\linewidth}{\centering\vspace{1.4cm}%
      \texttt{\detokenize{#1}}\ (captura de ecrã por inserir)\vspace{1.4cm}}}}}

\title{\textbf{Análise e Redução do Volume de Tráfego do FACEIT (CS2): Captura Global, Classificação, Filtragem e Agregação em Fluxos}}
\author{\textbf{GRUPO 6}\\[1.2ex]
Ivo Costa Sousa -- \texttt{pg63976}\\
João Afonso Almeida Sousa -- \texttt{pg62405}\\
João Carlos Teixeira Neiva -- \texttt{pg63978}}
\date{Novembro de 2026}

\begin{document}
\maketitle

\begin{abstract}
Este relatório, desenvolvido para a Unidade Curricular de Qualidade de Serviço na Internet, analisa o tráfego gerado pela
plataforma FACEIT durante uma partida de CS2. A partir de uma captura global com o Wireshark (326\,122 pacotes, 206,9~MB,
cerca de 30~minutos), o serviço é classificado num plano de jogo (um único fluxo UDP, 79\,\% dos bytes, sem marcação DSCP e com
um pacote a cada 15,6~ms) e num plano de controlo (DNS, TLS e QUIC, cerca de 1\,\% do volume). São estudados os critérios de
filtragem mais eficazes e quantificada a redução de volume obtida por filtragem (21,3\,\%), truncagem de pacotes (87,5\,\%),
\emph{sampling} combinado com truncagem (98,9\,\%) e agregação em fluxos com o Argus e com o \texttt{tshark} (99,96\,\% e 99,97\,\%).
\end{abstract}

%============================================================
\section{Parte 1: Objetivos e Metodologia}

\subsection{Objetivos}
Este trabalho consolida as etapas estudadas sobre captura e análise de tráfego, tendo como serviço de estudo o FACEIT,
plataforma competitiva de partidas de CS2 que exige a execução de um cliente de \emph{Anti-Cheat} próprio~\cite{faceitac}.
Os objetivos definidos foram:
\begin{itemize}
  \item \textbf{Captura global:} gerar tráfego real do serviço e capturar todo o tráfego da máquina com o Wireshark~\cite{wsug}.
  \item \textbf{Classificação e filtragem:} estudar os critérios mais eficazes, tanto em filtros de captura como em filtros
        de visualização.
  \item \textbf{Redução de volume:} quantificar a percentagem de otimização face à captura global por filtragem, truncagem de
        pacotes e \emph{sampling}.
  \item \textbf{Agregação em fluxos:} explorar ferramentas de agregação de pacotes em fluxos a partir de ficheiros \texttt{pcap}.
\end{itemize}

\subsection{Captura}
O Wireshark correu na própria máquina de jogo (\texttt{192.168.1.82}), ligada por cabo, \textbf{sem filtro de captura e com
pacotes inteiros}, durante 1\,818,6~s (cerca de 30~minutos) e com o FACEIT Anti-Cheat ativo. Antes da captura final realizou-se
um ensaio num servidor normal de CS2, fora do FACEIT, para validar o procedimento.

\subsection{Estratégia de Análise}
Optou-se por uma única captura global, convertida para \texttt{pcap} simples (\texttt{editcap -F pcap}) para servir de
\emph{baseline}. Todos os cenários reduzidos foram \textbf{derivados offline desse mesmo ficheiro}, garantindo que a comparação
incide exatamente sobre o mesmo tráfego. Recorreu-se ao Wireshark/\texttt{tshark}, \texttt{editcap}, \texttt{capinfos},
\texttt{tcpdump} e ao Argus~\cite{argus}. A percentagem de otimização foi calculada, em bytes (tamanho do ficheiro) e em número
de pacotes, por:
\[
  \text{redução} = \Bigl(1 - \frac{\text{tamanho reduzido}}{\text{tamanho global}}\Bigr)\times 100 .
\]
O tamanho do ficheiro inclui 16~B de cabeçalho por pacote no formato \texttt{pcap}. O campo \emph{Data size} do \texttt{capinfos}
não reflete a truncagem, pelo que não foi utilizado.

%============================================================
\section{Parte 2: Caracterização do Tráfego}

\subsection{Visão Geral}
A captura global contém 326\,122 pacotes e 201,6~MB de dados (212,6~MB em \texttt{pcapng} e 206,9~MB em \texttt{pcap}), com uma
média de 618~B por pacote e 179 pacotes/s. A composição por classe encontra-se na Tabela~\ref{tab:comp}; as classes
sobrepõem-se parcialmente, pelo que não devem ser somadas.

\begin{table}[htbp]
\centering
\caption{Composição do Tráfego Global}\label{tab:comp}
\begin{tabularx}{\linewidth}{@{}>{\bfseries\raggedright\arraybackslash}Xrrrr@{}}
\toprule
\textbf{Classe} & \textbf{Pacotes} & \textbf{Bytes} & \textbf{\% Pacotes} & \textbf{\% Bytes}\\
\midrule
Servidor de jogo (\texttt{141.95.32.214})     & 202\,086 & 159\,679\,469 & 62,0 & 79,2\\
UDP \texttt{:19336} (\texttt{104.29.146.69})  &  85\,785 &  17\,938\,347 & 26,3 &  8,9\\
TCP (todo)                                     &  19\,394 &  12\,376\,258 &  5,9 &  6,1\\
QUIC (todo)                                    &   7\,226 &   9\,622\,529 &  2,2 &  4,8\\
mDNS / SSDP / DNS                              &   4\,576 &   1\,307\,245 &  1,4 &  0,6\\
Controlo FACEIT (\texttt{193.41.200.101/102})  &   3\,995 &   3\,338\,873 &  1,2 &  1,7\\
\bottomrule
\end{tabularx}
\end{table}

A primeira linha conta todos os pacotes com o IP do servidor (202\,086), incluindo 80 que não pertencem à conversa UDP principal
(202\,006 pacotes), que é o valor usado no resto do relatório.

\subsection{Plano de Jogo}
A conversa UDP \texttt{192.168.1.82:50088}~$\leftrightarrow$~\texttt{141.95.32.214:27031} domina a captura
(Figura~\ref{fig:udp}). Inicia-se aos 293,8~s, dura 1\,520,5~s e totaliza 202\,006 pacotes e 159,7~MB. Os seus principais
indicadores são:
\begin{itemize}
  \item \textbf{Sentido cliente $\rightarrow$ servidor:} 96\,786 pacotes (70,7~MB), com média de 730~B por pacote.
  \item \textbf{Sentido servidor $\rightarrow$ cliente:} 105\,220 pacotes (89,0~MB), com média de 846~B por pacote.
  \item \textbf{Débito:} praticamente constante, com cerca de 133 pacotes/s e 0,84~Mbit/s (aproximadamente 40\,000 pacotes por
        cada 300~s).
  \item \textbf{Tamanho dos pacotes:} concentrado entre 640 e 1\,279~B (69,3\,\% dos pacotes), com média de 790~B.
\end{itemize}

\subsection{Características de QoS do Fluxo de Jogo}
Como a Unidade Curricular se centra na Qualidade de Serviço, caracterizou-se o fluxo de jogo do ponto de vista da marcação e
da regularidade temporal, medidas no ponto de captura (a máquina de jogo). Os 202\,006 pacotes têm \textbf{DSCP~0} nos dois
sentidos, ou seja, o tráfego é tratado como \emph{best effort}, sem qualquer marcação de QoS. A Tabela~\ref{tab:qos} resume o
tempo entre pacotes consecutivos de cada sentido.

\begin{table}[htbp]
\centering
\caption{Tempo entre Pacotes Consecutivos do Fluxo de Jogo (ms)}\label{tab:qos}
\begin{tabularx}{\linewidth}{@{}>{\bfseries\raggedright\arraybackslash}Xrrrrrrr@{}}
\toprule
\textbf{Sentido} & \textbf{Pacotes} & \textbf{Média} & \textbf{Mediana} & \textbf{p95} & \textbf{p99} & \textbf{Máximo} & \textbf{Desvio-padrão}\\
\midrule
Servidor $\rightarrow$ cliente & 105\,220 & 14,45 & 15,60 & 16,05 & 16,52 & 1\,000,4 & 7,58\\
Cliente $\rightarrow$ servidor &  96\,786 & 15,71 & 15,64 & 16,65 & 17,11 &   518,6 & 4,02\\
\bottomrule
\end{tabularx}
\end{table}

As medianas de cerca de 15,6~ms correspondem a aproximadamente 64 pacotes/s por sentido, e os percentis 95 e 99 ficam muito
próximos da mediana, o que mostra um fluxo muito regular. No sentido servidor $\rightarrow$ cliente, a média abaixo da mediana
e o desvio-padrão maior sugerem pacotes que chegam em rajadas, e os máximos de 1\,000~ms e 519~ms correspondem a interrupções
isoladas. Trata-se, portanto, de um fluxo de baixo débito mas sensível à latência e ao \emph{jitter}.

\subsection{Plano de Controlo}
O DNS e o SNI (em TLS e QUIC) permitem identificar inequivocamente o serviço. Foram observados os nomes
\nolinkurl{www.faceit.com}, \nolinkurl{edge.faceit.com}, \nolinkurl{chat.faceit.com}, \nolinkurl{faceit.com},
\nolinkurl{cdn-frontend.faceit-cdn.net}, \nolinkurl{distribution.faceit-cdn.net} e \nolinkurl{ac-client-ws.faceit.com}, todos
resolvidos para \texttt{193.41.200.101} e \texttt{193.41.200.102} (Figuras~\ref{fig:sni} e~\ref{fig:dns}), a que se junta
\nolinkurl{analytics.faceitanalytics.com} (Cloudflare). O primeiro contacto ocorre aos 2,2~s (QUIC para
\nolinkurl{www.faceit.com}), surge uma nova rajada de pedidos entre os 225 e os 245~s, antes de o fluxo de jogo começar, e um
último contacto a \nolinkurl{ac-client-ws.faceit.com} aos 1\,804,6~s. Este plano é reduzido: 3\,995 pacotes e 3,3~MB.

\subsection{Ruído}
Os fluxos UDP para \texttt{104.29.146.69:19336} (um novo fluxo por minuto, 85\,785 pacotes) são independentes do jogo, uma vez
que surgiam também na captura de ensaio; a sua origem não foi identificada. Há ainda uma ligação TCP de longa duração
(1\,818~s) para \texttt{185.25.182.20:27021}, provavelmente do cliente Steam. O restante tráfego corresponde à rede local (mDNS,
SSDP, \emph{broadcasts}) e a outras ligações HTTPS/QUIC.

\begin{figure}[htbp]
\centering
\figecra{fig_udp.png}
\caption{\emph{Conversations} $\rightarrow$ UDP, ordenado por pacotes: o fluxo para \texttt{141.95.32.214:27031} domina.}\label{fig:udp}
\end{figure}

\begin{figure}[htbp]
\centering
\figecra{fig_sni.png}
\caption{Filtro \texttt{tls.handshake.extensions\_\allowbreak server\_name contains "faceit"} (TLS e QUIC).}\label{fig:sni}
\end{figure}

\begin{figure}[htbp]
\centering
\figecra{fig_dns.png}
\caption{Filtro \texttt{dns.qry.\allowbreak name contains "faceit"}.}\label{fig:dns}
\end{figure}

%============================================================
\section{Parte 3: Critérios de Classificação e Filtragem}

\subsection{Filtros Estudados}
O filtro de captura (BPF)~\cite{pcapfilter} apenas tem acesso aos cabeçalhos (IP, protocolo, portas), pelo que critérios
baseados em nomes (DNS/SNI) só são possíveis como filtros de visualização. A eficácia de cada critério foi avaliada através da
\textbf{cobertura} (pacotes do jogo apanhados / pacotes do jogo) e da \textbf{precisão} (pacotes do jogo apanhados / pacotes
apanhados), tomando como referência a conversa do jogo (202\,006 pacotes). Os resultados constam da Tabela~\ref{tab:filtros}.

\begin{table}[htbp]
\centering
\caption{Critérios de Classificação e Respetiva Eficácia}\label{tab:filtros}
\footnotesize
\begin{tabularx}{\linewidth}{@{}>{\bfseries\raggedright\arraybackslash}p{2.7cm}YYrr@{}}
\toprule
\textbf{Critério} & \textbf{Filtro de Captura (BPF)} & \textbf{Filtro de Visualização} & \textbf{Pacotes} & \textbf{Precisão}\\
\midrule
Servidor (IP) & \texttt{udp and host 141.95.32.214} & \texttt{ip.addr==141.95.32.214 \&\& udp} & 202\,006 & 100\,\%\\
Qualquer protocolo com o IP & \texttt{host 141.95.32.214} & \texttt{ip.addr==141.95.32.214} & 202\,060 & 99,97\,\%\\
Gama de portas & \texttt{udp portrange 27000-27050} & \texttt{udp.port>=27000 \&\& udp.port<=27050} & 202\,040 & 99,98\,\%\\
Só UDP & \texttt{udp} & \texttt{udp} & 300\,571 & 67,2\,\%\\
Controlo (IP) & \texttt{host 193.41.200.101 or host 193.41.200.102} & \texttt{ip.addr==193.41.200.101 \textbar\textbar{} ip.addr==193.41.200.102} & 3\,995 & ---\\
Controlo (nome) & (não aplicável) & \texttt{dns.qry.\allowbreak name contains "faceit"}; \texttt{tls.handshake.extensions\_\allowbreak server\_name contains "faceit"} & --- & ---\\
Jogo + controlo (IPs) & \texttt{host} servidor \texttt{or host} controlo (2 IPs) & --- & 206\,055 & ---\\
Portas + controlo & \texttt{udp portrange 27000-27050 or host 193.41.200.101 or host 193.41.200.102} & --- & 206\,035 & ---\\
\bottomrule
\end{tabularx}
\end{table}

\subsection{Discussão}
Todos os filtros baseados em IP ou em portas atingem uma cobertura de 100\,\% do fluxo de jogo. A gama de portas apanha apenas
34 pacotes a mais do que o filtro por IP (precisão de 99,98\,\%) e tem a vantagem de não depender do IP do servidor, que muda de
partida para partida. Em contrapartida, o filtro apenas por protocolo (\texttt{udp}) revela-se demasiado abrangente, já que
33\,\% do que apanha não pertence ao jogo.

Conclui-se que a deteção do fluxo de jogo é mais eficaz através da combinação de três critérios: transporte UDP, gama de portas
\texttt{27xxx} e dominância por volume (\emph{Conversations} ordenado por bytes), sendo o serviço confirmado pelo DNS/SNI do plano
de controlo. Para uma captura ao vivo, em que o IP do servidor ainda não é conhecido, o filtro mais adequado é o último da
Tabela~\ref{tab:filtros}, que junta a gama de portas aos dois IPs do plano de controlo.

%============================================================
\section{Parte 4: Redução de Volume}

\subsection{Resultados}
A Tabela~\ref{tab:red} e a Figura~\ref{fig:red} apresentam os cenários avaliados. Todas as percentagens são cumulativas face ao
\emph{baseline} (A) de 326\,122 pacotes e 206\,864\,624~B.

\begin{figure}[!htbp]
\centering
\begin{tikzpicture}
\begin{axis}[
  xbar, width=0.74\linewidth, height=6.6cm,
  xmin=0, xmax=115, xlabel={Redução face à captura global (\% dos bytes)},
  symbolic y coords={B1 filtro IP, B3 filtro portas, B6 portas+controlo, C1 snaplen 128, C2 snaplen 96, C3 snaplen 64, C5 filtro recomendado, D1 sampling 1/10, D2 snaplen 96 + sampling, E1 Argus 5 s, E4 tshark conversas, E2 Argus 1/fluxo},
  ytick=data, y dir=reverse, enlarge y limits=0.06, bar width=7pt,
  yticklabel style={font=\footnotesize}, xmajorgrids,
  nodes near coords={\pgfmathprintnumber[use comma,fixed,fixed zerofill,precision=2]{\pgfplotspointmeta}},
  every node near coord/.append style={font=\scriptsize}
]
\addplot[fill=blue!55,draw=blue!70!black] coordinates {
  (21.25,B1 filtro IP) (21.23,B3 filtro portas) (19.58,B6 portas+controlo)
  (84.38,C1 snaplen 128) (87.50,C2 snaplen 96) (90.63,C3 snaplen 64) (87.27,C5 filtro recomendado)
  (92.08,D1 sampling 1/10) (98.91,D2 snaplen 96 + sampling)
  (99.51,E1 Argus 5 s) (99.96,E4 tshark conversas) (99.97,E2 Argus 1/fluxo)};
\end{axis}
\end{tikzpicture}
\caption{Redução de Bytes por Cenário}\label{fig:red}
\end{figure}

\begin{table}[htbp]
\centering
\caption{Redução de Volume face à Captura Global}\label{tab:red}
\footnotesize
\begin{tabularx}{\linewidth}{@{}>{\bfseries}lXrrrr@{}}
\toprule
 & \textbf{Cenário} & \textbf{Pacotes} & \textbf{Bytes} & \textbf{Red. Bytes} & \textbf{Red. Pacotes}\\
\midrule
A  & Captura global (\emph{baseline})                          & 326\,122 & 206\,864\,624 & 0\,\%     & 0\,\%\\
B1 & Filtro \texttt{udp and host} servidor                     & 202\,006 & 162\,905\,309 & 21,25\,\% & 38,06\,\%\\
B3 & Filtro gama de portas 27000--27050                        & 202\,040 & 162\,951\,481 & 21,23\,\% & 38,05\,\%\\
B4 & Filtro só \texttt{udp}                                    & 300\,571 & 193\,613\,964 &  6,41\,\% &  7,83\,\%\\
B5 & Filtro jogo + controlo FACEIT (IPs)                       & 206\,055 & 166\,312\,206 & 19,60\,\% & 36,82\,\%\\
B6 & Filtro portas + controlo FACEIT                           & 206\,035 & 166\,354\,274 & 19,58\,\% & 36,82\,\%\\
C1 & B1 + \emph{snaplen} 128                                   & 202\,006 &  32\,311\,732 & 84,38\,\% & 38,06\,\%\\
C2 & B1 + \emph{snaplen} 96                                    & 202\,006 &  25\,856\,584 & 87,50\,\% & 38,06\,\%\\
C3 & B1 + \emph{snaplen} 64                                    & 202\,006 &  19\,392\,644 & 90,63\,\% & 38,06\,\%\\
C4 & B5 + \emph{snaplen} 96                                    & 206\,055 &  26\,333\,772 & 87,27\,\% & 36,82\,\%\\
C5 & B6 + \emph{snaplen} 96 (filtro recomendado)               & 206\,035 &  26\,333\,156 & 87,27\,\% & 36,82\,\%\\
D1 & B1 + \emph{sampling} 1/10                                 &  20\,200 &  16\,376\,628 & 92,08\,\% & 93,81\,\%\\
D2 & B1 + \emph{snaplen} 96 + \emph{sampling} 1/10             &  20\,200 &   2\,262\,424 & 98,91\,\% & 93,81\,\%\\
E1 & Fluxos Argus, global (registos de 5~s)                    & --- &   1\,004\,872 & 99,51\,\% & ---\\
E2 & Fluxos Argus, global (1 registo/fluxo, CSV)               & --- &      64\,839 & 99,97\,\% & ---\\
E3 & Fluxos Argus, só servidor (CSV, 1 fluxo)                  & --- &         229 & $\approx$100\,\% & ---\\
E4 & Conversas \texttt{tshark}, global (texto, UDP+TCP)        & --- &      89\,824 & 99,96\,\% & ---\\
\bottomrule
\end{tabularx}
\end{table}


\subsection{Filtragem}
Filtrar o fluxo de jogo reduz apenas 21,25\,\% dos bytes, mas 38,06\,\% dos pacotes. Esta assimetria explica-se pelo facto de o
jogo já representar 79\,\% dos bytes capturados, enquanto o ruído é composto por muitos pacotes pequenos: os 85\,785 pacotes de
\texttt{:19336} correspondem a 26\,\% dos pacotes, mas a apenas 9\,\% dos bytes. Incluir o plano de controlo no filtro tem um
custo reduzido (19,58\,\% a 19,60\,\% contra 21,25\,\%) e descreve o serviço de forma mais completa.

\subsection{Truncagem de Pacotes (Snaplen)}
Não capturar o pacote inteiro é a técnica com maior impacto individual. Os cabeçalhos Ethernet, IPv4 e UDP ocupam 42~B, pelo
que, com \emph{snaplen} 96, se preservam os cabeçalhos e mais 54~B de \emph{payload}. Com pacotes de 790~B em média, a redução
total sobe de 21,25\,\% para 87,50\,\% (84,38\,\% com 128~B e 90,63\,\% com 64~B). O custo é a perda do restante \emph{payload}.

\subsection{Sampling}
Guardar 1 pacote em cada 10 (por ordem de pacote) proporciona uma redução de 92,08\,\% sem truncagem e de 98,91\,\% com
\emph{snaplen} 96 (2,26~MB, dos quais cerca de 0,32~MB correspondem a cabeçalhos de pacote do formato \texttt{pcap}). As
estatísticas de volume mantêm-se fiáveis: extrapolando $\times 10$, o erro é de $-0{,}003\,\%$ nos pacotes e de $+0{,}54\,\%$
nos bytes. Perdem-se, contudo, as sequências de pacotes consecutivos.

\subsection{Impacto nas Métricas de QoS}
As técnicas não são equivalentes para quem quer medir QoS. A filtragem e o \emph{snaplen} 96 mantêm todos os pacotes, as marcas
temporais e o cabeçalho IP (onde está o DSCP), pelo que ainda permitem calcular o tempo entre pacotes e o \emph{jitter} da
Tabela~\ref{tab:qos}. O \emph{sampling} 1/10 preserva volumes e distribuições de tamanhos, mas os intervalos medidos entre
pacotes amostrados deixam de corresponder aos reais. A agregação em fluxos conserva o débito médio e a duração, mas perde o
\emph{jitter}.

%============================================================
\section{Parte 5: Agregação de Pacotes em Fluxos}

\subsection{Ferramentas e Processo}
Exploraram-se duas ferramentas sobre o mesmo \texttt{A\_global.pcap}. O \textbf{Argus} (\texttt{argus -r f.pcap -w f.argus})
emite registos de estado de cerca de 5~s por fluxo, e o \texttt{racluster} agrupa-os num único registo por fluxo (5-tuplo),
exportado para CSV através do \texttt{ra}. O \textbf{\texttt{tshark}} produz, com \texttt{-z conv,udp -z conv,tcp}, uma lista
de conversas UDP e TCP em texto.

\subsection{Resultados}
A Tabela~\ref{tab:fluxos} compara as saídas obtidas a partir do \emph{baseline} de 206,9~MB.

\begin{table}[htbp]
\centering
\caption{Agregação em Fluxos: Argus e tshark}\label{tab:fluxos}
\footnotesize
\begin{tabularx}{\linewidth}{@{}>{\bfseries}Yrrrr@{}}
\toprule
\textbf{Saída} & \textbf{Registos} & \textbf{Tamanho (B)} & \textbf{Redução} & \textbf{Protocolos}\\
\midrule
Argus, registos de 5~s (binário)          & ---  & 1\,004\,872 & 99,51\,\% & todos\\
Argus, 1 registo por fluxo (binário)      & 797  &   122\,656 & 99,94\,\% & todos\\
Argus, 1 registo por fluxo (CSV)          & 797  &    64\,839 & 99,97\,\% & todos\\
\texttt{tshark} conversas (texto)         & 612  &    89\,824 & 99,96\,\% & UDP e TCP\\
\bottomrule
\end{tabularx}
\end{table}

As duas ferramentas chegam a ordens de grandeza semelhantes (cerca de 65~kB a 90~kB, mais de 99,9\,\%), mas diferem no que
registam:
\begin{itemize}
  \item \textbf{Cobertura de protocolos:} o \texttt{tshark} só lista UDP e TCP (447 e 165 conversas, tal como mostram os
        separadores das \emph{Conversations}). O Argus regista também ICMPv6, ARP, IGMP, ICMP e LLDP, daí os 797 fluxos
        (mais um registo de gestão do Argus), dos quais cerca de 200 não são UDP nem TCP.
  \item \textbf{Fluxos UDP e TCP:} no Argus há 448 fluxos UDP e 149 TCP, contra 447 e 165 no \texttt{tshark}. A diferença nos
        TCP não foi investigada; provavelmente deve-se a critérios distintos de delimitação das conversas.
  \item \textbf{Precisão dos totais:} o \texttt{tshark} em texto arredonda os bytes (por exemplo, \emph{159~MB} para o fluxo de
        jogo), enquanto o Argus conserva os valores exatos.
\end{itemize}

Os totais do fluxo de jogo no Argus (70\,661\,525~B do cliente e 89\,011\,664~B do servidor) somam 159\,673\,189~B, exatamente
o volume de dados do \texttt{pcap} filtrado. Aplicado ao filtro do servidor, o fluxo de jogo fica reduzido a um único registo de
229~B no CSV. A agregação conserva o 5-tuplo, o instante de início, a duração e os pacotes e bytes por sentido, mas perde a
temporização e o tamanho de cada pacote, bem como o \emph{payload}. Os registos de 5~s do Argus conservam ainda a evolução
temporal do fluxo.

%============================================================
\section{Parte 6: Conclusões}

\subsection{Síntese}
\begin{itemize}
  \item \textbf{Estrutura do serviço:} o tráfego do FACEIT divide-se num plano de jogo (UDP \texttt{141.95.32.214:27031}, 79\,\%
        dos bytes) e num plano de controlo (DNS/TLS/QUIC para \texttt{193.41.200.101/102}, cerca de 1\,\%).
  \item \textbf{Classificação:} os critérios mais eficazes para detetar o fluxo de jogo são o transporte UDP, a gama de portas
        27000--27050 e a dominância por volume; o DNS/SNI confirma o serviço. O filtro apenas por protocolo é demasiado abrangente.
  \item \textbf{QoS do fluxo de jogo:} DSCP~0 em todos os pacotes e um pacote a cada 15,6~ms em cada sentido, com pouco desvio
        em torno desse valor.
  \item \textbf{Redução de volume:} a filtragem poupa pouco em bytes (21\,\%), dado que o jogo domina a captura. As maiores
        poupanças resultam do \emph{snaplen} (87,5\,\%), do \emph{sampling} (98,9\,\% em conjunto) e da agregação em fluxos
        (99,5\,\% a 99,97\,\%).
  \item \textbf{Escolha da técnica:} depende da utilização pretendida: cabeçalhos e tempos por pacote $\Rightarrow$ filtro e
        \emph{snaplen} 96; estatísticas de volume $\Rightarrow$ \emph{sampling}; contabilização e retenção prolongada
        $\Rightarrow$ fluxos.
\end{itemize}

\subsection{Configuração de Captura Recomendada}
Para uma captura ao vivo deste serviço recomenda-se o filtro de portas e IPs de controlo com \emph{snaplen} 96, por exemplo
com o \texttt{dumpcap}:
\begin{lstlisting}
dumpcap -i <n> -f "udp portrange 27000-27050 or host 193.41.200.101 or host 193.41.200.102" -s 96 -w captura.pcapng
\end{lstlisting}
Aplicada offline à captura global, esta configuração (C5) guarda 206\,035 pacotes em 26,3~MB, uma redução de 87,27\,\% dos
bytes, e mantém os cabeçalhos, as marcas temporais e o DNS/SNI dos pedidos de controlo apenas até aos primeiros 96~B de cada
pacote. Não foi validada numa captura ao vivo.

\subsection{Limitações}
Foi analisada uma única captura de uma única partida, e o IP e a porta do servidor mudam entre partidas. Os cenários foram
derivados offline da captura global, e não capturados ao vivo. O \emph{sampling} aplicado é determinístico (1 em 10 por ordem de
pacote) e as percentagens em ficheiro incluem o cabeçalho de 16~B por pacote do formato \texttt{pcap}. O tempo entre pacotes
foi medido no ponto de captura e inclui, por isso, o efeito da rede e do próprio sistema de captura. A agregação em fluxos
usou apenas o Argus e o \texttt{tshark}.

%============================================================
{\small
\begin{thebibliography}{9}
\bibitem{wsug} Wireshark Foundation, \emph{Wireshark User's Guide}. \url{https://www.wireshark.org/docs/wsug_html_chunked/}
\bibitem{pcapfilter} The Tcpdump Group, \emph{pcap-filter(7) --- packet filter syntax}. \url{https://www.tcpdump.org/manpages/pcap-filter.7.html}
\bibitem{argus} QoSient, \emph{Argus --- Auditing Network Activity}. \url{https://openargus.org}
\bibitem{faceitac} FACEIT, \emph{What is FACEIT Anti-cheat and how does it work?}
\url{https://support.faceit.com/hc/en-us/articles/9394666828188-What-is-FACEIT-Anti-cheat-and-how-does-it-work}
\end{thebibliography}
}

%============================================================
\appendix
\lstset{basicstyle=\ttfamily\scriptsize}
\section{Comandos Utilizados}
\begin{lstlisting}
# Baseline
capinfos faceit.pcapng
editcap -F pcap faceit.pcapng A_global.pcap

# Caracterizacao
tshark -r A_global.pcap -q -z io,phs
tshark -r A_global.pcap -q -z conv,udp
tshark -r A_global.pcap -q -z plen,tree,"ip.addr==141.95.32.214"
tshark -r A_global.pcap -Y 'dns.flags.response==1 && dns.qry.name contains "faceit"' -T fields -e frame.time_relative -e dns.qry.name -e dns.a
tshark -r A_global.pcap -Y 'tls.handshake.extensions_server_name contains "faceit"' -T fields -e frame.time_relative -e ip.dst -e tls.handshake.extensions_server_name

# QoS do fluxo de jogo (DSCP e instantes; os intervalos foram calculados em Python)
tshark -r B1_udp_host.pcap -T fields -e ip.dsfield.dscp | sort | uniq -c
tshark -r B1_udp_host.pcap -Y "ip.src==141.95.32.214" -T fields -e frame.time_epoch

# Filtros (BPF, aplicados offline)
tcpdump -r A_global.pcap -w B1_udp_host.pcap  "udp and host 141.95.32.214"
tcpdump -r A_global.pcap -w B3_portrange.pcap "udp portrange 27000-27050"
tcpdump -r A_global.pcap -w B4_udp.pcap       "udp"
tcpdump -r A_global.pcap -w B5_game_ctrl.pcap "host 141.95.32.214 or host 193.41.200.101 or host 193.41.200.102"
tcpdump -r A_global.pcap -w B6_portrange_ctrl.pcap "udp portrange 27000-27050 or host 193.41.200.101 or host 193.41.200.102"

# Snaplen e sampling
editcap -s 96 B1_udp_host.pcap C2_snaplen96.pcap
editcap -s 96 B6_portrange_ctrl.pcap C5_portrange_ctrl_snaplen96.pcap
tshark -r C2_snaplen96.pcap -Y "frame.number % 10 == 0" -F pcap -w D2_snaplen_sampling.pcap

# Fluxos (Argus)
argus -r A_global.pcap -w E1_global.argus
racluster -r E1_global.argus -m proto saddr sport daddr dport -w E2_global_cluster.argus
ra -r E2_global_cluster.argus -n -c , -s stime dur proto saddr sport daddr dport spkts dpkts sbytes dbytes > E2_global_flows.csv

# Fluxos (tshark) e medicao
tshark -r A_global.pcap -q -z conv,udp -z conv,tcp > flows_tshark.txt
capinfos -T -m -M -c -s *.pcap
\end{lstlisting}

\end{document}